%% kapa: exam=1 round=32 year=2021 period=1 form=A subject=real_estate subject_ko=부동산학원론 seq=6 number=86
%% source: Q-Net 감정평가사 기출문제 (cst00309) artlSeq=5208397 file=extracted/1008/5208397/A형/2021년도 제32회 감정평가사 1차 1교시 A형.pdf
%% conversion: kordoc 4.24.51 -> tools/gen1.py (auto); review=proofread-sonnet; flags=none
\begin{question}{86}
도시공간구조이론에 관한 설명으로 옳지 않은 것은?
\begin{choices}
  \item 동심원이론은 도시공간구조의 형성을 침입, 경쟁, 천이 과정으로 설명하였다.
  \item 동심원이론에 따르면 중심지에서 멀어질수록 지대 및 인구밀도가 낮아진다.
  \item 선형이론에서의 점이지대는 중심업무지구에 직장 및 생활터전이 있어 중심업무지구에 근접하여 거주하는 지대를 말한다.
  \item 선형이론에 따르면 도시공간구조의 성장 및 분화가 주요 교통노선을 따라 부채꼴 모양으로 확대된다.
  \item 다핵심이론에 따르면 하나의 중심이 아니라 몇 개의 분리된 중심이 점진적으로 통합됨에 따라 전체적인 도시공간구조가 형성된다.
\end{choices}
\end{question}
\answer{3}
